% Part: computability
% Chapter: computability-theory
% Section: prop-reduce

\documentclass[../../../include/open-logic-section]{subfiles}

\begin{document}

\olfileid{cmp}{thy}{ppr}
\olsection{ન્યૂનીકરણક્ષમતાના ગુણધર્મો}

$A$ને~$B$માં ઘટાડી શકાય ત્યારે $A \leq_m B$ લખીએ છીએ. આ
સંકેત સૂચવે છે કે $A \leq_m B$ હોય તો $A$ એ~$B$ કરતાં
``વધુ કઠિન નથી'' અને $B$ એ~$A$ જેટલો અથવા વધુ કઠિન છે.
સંખ્યાઓ પરના સામાન્ય $\le$ની જેમ $\le_m$ ગણો (અથવા
નિર્ણય સમસ્યાઓ)ને તેમની જટિલતા મુજબ ક્રમિત કરે છે.
આગળના બે વિધાનો આ સમજને ટેકો આપે છે. પહેલું કહે છે કે
$\le_m$ સંક્રામક છે.

\begin{prop}
\ollabel{prop:trans-red}
$A \leq_m B$ અને $B \leq_m C$ હોય, તો $A \leq_m C$.
\end{prop}

\begin{proof}
$A$નું~$B$માં ન્યૂનીકરણ અને $B$નું $C$માં ન્યૂનીકરણ
સંયોજવાથી $A$નું~$C$માં ન્યૂનીકરણ મળે છે.
\end{proof}

\begin{prob}
જો $f$ અને~$g$ અનુક્રમે $A$નું~$B$માં અને $B$નું~$C$માં
બહુ-એક ન્યૂનીકરણ હોય, તો $\comp{f}{g}$ એ $A$નું~$C$માં
બહુ-એક ન્યૂનીકરણ છે એવું બતાવી \olref{prop:trans-red}
સાબિત કરો.
\end{prob}

\begin{prop}
\ollabel{prop:reduce}
$A$ અને $B$ કોઈ પણ ગણો હોય અને $A \leq_m B$ માનો.
\begin{enumerate}
\item $B$ સંગણકીય રીતે પરિગણનીય હોય, તો $A$ પણ છે.
\item $B$ સંગણનીય હોય, તો $A$ પણ છે.
\end{enumerate}
\end{prop}

\begin{proof}
$f$, $A$નું~$B$માં બહુ-એક ન્યૂનીકરણ હોય. પહેલા દાવા
માટે તપાસો કે $B$ આંશિક વિધેય~$g$નો પ્રદેશ હોય, તો
$A$ એ~$\comp{f}{g}$નો પ્રદેશ છે:
\begin{align*}
x \in A & \text{ ત્યારે અને માત્ર ત્યારે } f(x) \in B \\
& \text{ ત્યારે અને માત્ર ત્યારે }  g(f(x)) \fdefined.
\end{align*}

બીજા દાવા માટે યાદ કરો કે $B$ સંગણનીય હોય, તો $B$
અને~$\Complement{B}$ બંને સંગણકીય રીતે પરિગણનીય છે
(\olref[cmp]{thm:ce-comp}). $f$, $\Complement{A}$નું
$\Complement{B}$માં પણ બહુ-એક ન્યૂનીકરણ છે તે તપાસવું
મુશ્કેલ નથી. તેથી સાબિતીના પહેલા ભાગથી $A$
અને~$\Complement{A}$ બંને !!{computably
enumerable} છે. આથી \olref[cmp]{thm:ce-comp} મુજબ $A$ પણ
સંગણનીય છે. (વૈકલ્પિક રીતે, $\Char{A} =
\comp{f}{\Char{B}}$ છે તે તપાસી શકાય; તેથી $\Char{B}$
સંગણનીય હોય, તો~$\Char{A}$ પણ છે.)
\end{proof}

\begin{prob}
માનો કે $f$, $A$નું~$B$માં બહુ-એક ન્યૂનીકરણ છે. બતાવો કે
$f$, $\Complement{A}$નું $\Complement{B}$માં પણ બહુ-એક
ન્યૂનીકરણ છે.
\end{prob}

\begin{prob}
$f\colon A \to B$ બહુ-એક ન્યૂનીકરણ હોય, તો
$\Char{A} = \comp{f}{\Char{B}}$ છે એવું બતાવો.
\end{prob}

\begin{digress}
\emph{ટ્યુરિંગ ન્યૂનીકરણક્ષમતા} કહેવાતી ન્યૂનીકરણક્ષમતાની
વધુ સામાન્ય કલ્પના બીજા સંદર્ભોમાં, ખાસ કરીને
અનિર્ણેયતાનાં પરિણામો સાબિત કરવા, ઉપયોગી છે. નોંધો કે
\olref[cmp]{cor:comp-k} મુજબ $K_0$નો પૂરક~$K_0$માં
ઘટાડી શકાતો નથી, કારણ કે તે સંગણકીય રીતે પરિગણનીય નથી.
પરંતુ સહજ રીતે, $K_0$ અંગેના પ્રશ્નોના જવાબો ખબર હોય, તો
તેના પૂરક અંગેના પ્રશ્નોના જવાબો પણ ખબર પડે. ગણ $A$ને~$B$માં
ટ્યુરિંગ-ન્યૂનીકરણક્ષમ ત્યારે કહીએ જ્યારે~$B$ વિશે પ્રશ્નો
પૂછી શકતી સંગણનીય પ્રક્રિયા વડે~$A$ અંગેના પ્રશ્નોના જવાબો
નક્કી કરી શકાય. આ બહુ-એક ન્યૂનીકરણક્ષમતા કરતાં ઉદાર છે:
તેમાં (1)~$B$ વિશે માત્ર એક જ પ્રશ્ન પૂછી શકાય છે અને
(2)~``હા'' જવાબનો $A$ વિશેના પ્રશ્ન માટે ``હા''માં,
અને ``ના''નો ``ના''માં અનુવાદ થવો જોઈએ. તેમ છતાં,
$A$ને~$B$માં ટ્યુરિંગ રીતે ઘટાડી શકાય અને $B$ સંગણનીય
હોય, તો $A$ પણ સંગણનીય છે; જોકે, આપણે જોયું તેમ,
સંગણકીય પરિગણનીયતા માટે સમાન વિધાન સાચું નથી.

આપણે ચર્ચેલી ન્યૂનીકરણક્ષમતાની જુદી કલ્પનાઓ અને તેમના
ભેદો વિશે વિચારો. આ પ્રકરણમાં, જોકે, માત્ર બહુ-એક
ન્યૂનીકરણક્ષમતા જ લઈશું. ઉપરના ફકરામાં ચર્ચેલા બંને પ્રકારનાં
ન્યૂનીકરણને સંગણકીય જટિલતામાં પણ અનુરૂપ પ્રકારો છે, જેમાં
ટ્યુરિંગ મશીનો બહુપદી સમયમાં ચાલે એવી વધારાની શરત છે:
બહુ-એક ન્યૂનીકરણક્ષમતાના જટિલતા-રૂપને \emph{કાર્પ
ન્યૂનીકરણક્ષમતા} અને ટ્યુરિંગ ન્યૂનીકરણક્ષમતાના રૂપને
\emph{કુક ન્યૂનીકરણક્ષમતા} કહે છે.
\end{digress}

\end{document}
